% Chapter 9: Discussion

\section{Introduction}

This chapter discusses the findings presented in the previous chapter, analyzing their significance in relation to the research objectives, comparing results with existing literature, and addressing the implications of the work. The discussion also examines limitations and provides insights into the practical applications of the Lesaka AI system.

\section{Interpretation of Results}

\subsection{Database Performance}

The excellent database performance results (sub-20ms for most operations) demonstrate that the 3NF schema, while theoretically more complex, does not introduce significant performance overhead for the expected data volumes in agricultural IoT contexts. This aligns with \cite{date2003database} findings that properly normalized databases can maintain performance while ensuring data integrity.

The 100\% success rate for referential integrity validation confirms that the foreign key constraints effectively prevent orphaned records, a critical requirement for maintaining data quality in telemetry systems.

\subsection{RBAC Effectiveness}

The RBAC implementation showed 100\% success in denying broker access to GPS coordinates, successfully implementing privacy-by-design principles. The 80\% success rate for broker access attempts reflects intentional restriction of sensitive data, not system failure. This demonstrates that the system effectively balances data utility with privacy protection, addressing the concerns raised by \cite{cavoukian2011privacy}.

\subsection{Multi-Agent System Performance}

The 97.7-98.6\% accuracy in agent routing demonstrates that the graph-based orchestration effectively categorizes telemetry data based on contextual factors. The 100\% self-healing recovery rate indicates that the fallback mechanism successfully prevents system deadlocks, addressing the resilience requirements for network-constrained rural environments.

\section{Comparison with Literature}

\subsection{IoT in Agriculture}

Compared to existing livestock monitoring systems \cite{riquelme2018smonitor}, the Lesaka AI system provides superior data governance through 3NF validation and RBAC, features notably absent in many existing solutions. The integration of intelligent agent orchestration represents a significant advancement over basic monitoring systems.

\subsection{Multi-Agent Systems}

The graph-based orchestration approach differs from traditional hierarchical agent systems \cite{wooldridge2009introduction} by enabling dynamic routing based on contextual factors. The self-healing mechanism addresses a gap identified in \cite{ghosh2007self} regarding fault tolerance in distributed systems.

\subsection{Data Governance}

The implementation of 3NF validation in an agricultural IoT context represents a novel application of database normalization principles, which are typically discussed in theoretical contexts rather than practical IoT deployments. The privacy-by-design implementation goes beyond the basic security measures found in most agricultural systems.

\section{Implications}

\subsection{Academic Implications}

The project demonstrates the successful integration of INFS and COMP module requirements in a single cohesive system, providing a model for interdisciplinary final year projects. The student-like code characteristics, including honest reflection on limitations and temporary fixes, align with the COMP402 assessment criteria that value understanding over polished output.

\subsection{Practical Implications}

For Botswana's agricultural sector, the system provides:

\begin{itemize}
    \item A scalable framework for livestock management
    \item Privacy protection for farmer data
    \item Intelligent diagnostic capabilities
    \item Resilient operation in network-constrained environments
\end{itemize}

The Botswana-specific customizations (Pula currency, regional districts, local pricing) demonstrate the importance of context-aware system design for technology adoption in developing regions.

\subsection{Technical Implications}

The project demonstrates that:

\begin{itemize}
    \item 3NF database design can be effectively implemented in IoT edge computing
    \item Multi-agent systems can provide self-healing capabilities
    \item RBAC can effectively implement privacy-by-design principles
    \item React-based frontends can be successfully integrated with Python backends
\end{itemize}

\section{Limitations}

\subsection{Technical Limitations}

\begin{itemize}
    \item \textbf{Async Implementation:} Full async/await implementation was deferred due to complexity and time constraints
    \item \textbf{Thread Safety:} Comprehensive thread safety testing was partial due to limited access to load testing environments
    \item \textbf{React Build:} npm unavailability prevented complete React build testing
    \item \textbf{Real-World Data:} Testing relied on synthetic data rather than real-world sensor inputs
\end{itemize}

\subsection{Scope Limitations}

\begin{itemize}
    \item \textbf{Hardware Integration:} Actual sensor hardware integration was out of scope
    \item \textbf{Mobile Application:} Mobile app development was deferred to future work
    \item \textbf{Machine Learning:} ML model training and deployment was not implemented
    \item \textbf{Blockchain:} Blockchain integration for provenance tracking was deferred
\end{itemize}

\subsection{Methodological Limitations}

\begin{itemize}
    \item \textbf{Single Developer:} The project was developed by a single developer, limiting code review depth
    \item \item \textbf{Time Constraints:} Academic semester timeline limited comprehensive testing
    \item \textbf{User Testing:} Limited access to real farmers for user acceptance testing
\end{itemize}

\section{Design Decisions and Trade-offs}

\subsection{SQLite vs PostgreSQL}

\textbf{Decision:} SQLite was chosen over PostgreSQL for the database.

\textbf{Rationale:} SQLite requires no separate database server, making it ideal for edge deployment in rural areas with limited infrastructure.

\textbf{Trade-off:} PostgreSQL would provide better concurrency and advanced features, but at the cost of deployment complexity.

\subsection{Flask vs FastAPI}

\textbf{Decision:} Flask was chosen over FastAPI for the web framework.

\textbf{Rationale:} Flask's simplicity and extensive documentation made it suitable for rapid development within academic timeline constraints.

\textbf{Trade-off:} FastAPI would provide native async support and automatic API documentation, but required more learning time.

\subsection{React vs Vanilla JS}

\textbf{Decision:} React was chosen over vanilla JavaScript for the frontend.

\textbf{Rationale:} React's component-based architecture and ecosystem enabled rapid development of complex UIs from Figma designs.

\textbf{Trade-off:} Vanilla JS would have a smaller bundle size and simpler deployment, but would require more development time for complex features.

\section{AI Assistance Reflection}

\subsection{AI Usage}

AI assistance (Cascade) was used for:

\begin{itemize}
    \item Code generation for standard patterns (database connections, API endpoints)
    \item Documentation template creation
    \item LaTeX dissertation structure
    \item Testing framework setup
\end{itemize}

\subsection{AI Limitations}

AI limitations encountered:

\begin{itemize}
    \item Generated code sometimes lacked context-specific optimizations
    \item AI suggested over-engineered solutions inappropriate for academic timeline
    \item AI struggled with Botswana-specific context without explicit guidance
\end{itemize}

\subsection{Human Improvements}

Human improvements made to AI-generated code:

\begin{itemize}
    \item Added student-like comments and temporary fixes
    \item Simplified over-engineered solutions
    \item Added Botswana-specific customizations
    \item Implemented honest error handling and logging
\end{itemize}

\subsection{Understanding Verification}

Understanding was verified through:

\begin{itemize}
    \item Code review and explanation of each component
    \item Manual testing and debugging
    \item Documentation of design decisions
    \item Honest reflection on limitations in dissertation
\end{itemize}

\section{Lessons Learned}

\subsection{Technical Lessons}

\begin{itemize}
    \item 3NF database design requires careful planning but pays off in data quality
    \item Multi-agent systems add complexity but provide valuable flexibility
    \item Self-healing mechanisms are essential for resilient IoT systems
    \item RBAC implementation requires careful role definition and permission mapping
\end{itemize}

\subsection{Project Management Lessons}

\begin{itemize}
    \item Agile development with regular commits aids progress tracking
    \item Comprehensive documentation from the start saves time later
    \item Early testing prevents major refactoring later
    \item Honest reflection on limitations is valued over perfect code
\end{itemize}

\subsection{Academic Lessons}

\begin{itemize}
    \item Understanding and justification are more important than polished output
    \item Interdisciplinary projects require careful requirement alignment
    \item Student-like code characteristics demonstrate authentic development
    \item Comprehensive documentation supports successful viva defence
\end{itemize}

\section{Conclusion}

The discussion demonstrates that the Lesaka AI system successfully addresses the research objectives while providing valuable insights into the integration of data governance, multi-agent systems, and privacy protection in agricultural IoT contexts. The limitations identified provide clear directions for future work, and the honest reflection on AI assistance aligns with COMP402 assessment criteria that value understanding over artificial perfection.

The system represents a solid foundation for continued development and deployment in Botswana's agricultural sector, with potential for expansion to other regions and agricultural contexts.
